\documentclass[12pt]{article}
%\linespread{1.1}
% \bibliographystyle{plain}
\usepackage[2632449]{easymcm}
\problem{C}
\usepackage{palatino} 
\usepackage{pdfpages}
\usepackage{longtable}
\usepackage{tabu}
\usepackage{threeparttable}
\usepackage{listings}
\usepackage{paralist}
\usepackage{microtype}
\usepackage{appendix}
\usepackage{xeCJK}
\setCJKmainfont{SimSun}
\graphicspath{{img/}}
\let\itemize\compactitem
\let\enditemize\endcompactitem
\newcommand{\upcite}[1]{\textsuperscript{\textsuperscript{\cite{#1}}}}
\title{ 龚王博是我爹 } 
\renewcommand{\contest}{MCM/ICM}
\newcommand\specialsectioning{\setcounter{secnumdepth}{-2}}
\newcommand\specialappendix{\setcounter{appletdepth}{-2}}


\begin{document}


\begin{abstract}
    
% \begin{figure}[htbp]
% 	\centering
% 	\includegraphics[scale=0.12]{Food Web.jpg}
% 	\caption{Food Web}
% 	\label{Figure 3}
% \end{figure}

    \vspace{5pt}
    \textbf{Keywords}:

\end{abstract}

\maketitle % summary

\tableofcontents % contents
 
 

% Chapter 1: Introduction
\section{Introduction}

\subsection{Problem Background}


\subsection{Restatement of the Problem}

\begin{itemize}
   \item 
   \item 
   \item 
   \item 
   \item 
\end{itemize}


% \subsection{Our Work}
% \begin{figure}[h]
%     \begin{center}
%     \includegraphics[scale=0.16]{Our Work.jpg}
%     \caption{Our Work}
%     \end{center}
% \end{figure}


% Chapter 2: 
\section{}

\subsection{}
\begin{enumerate} %有序列表
   \item {\bf 11}
   \item {\bf 22}
   \item {\bf 33}
\end{enumerate}


\subsection{Notations}
Some important mathematical notations used in this paper are listed in Table 1. 
\begin{table}[htbp]
    \begin{center}
        \caption{Notations used in this paper}
        \begin{tabular}{c c}
            \toprule[2pt]
            {\centering Symbol}
            &{\centering Description }\\
            \midrule
            $E$ & entropy \\
            \bottomrule[2pt]
        \end{tabular}
        \begin{tablenotes}
            \footnotesize
            \item[*] *There are some variables that are not listed here and 
            will be discussed in detail in each section. %此处加入注释*信息
        \end{tablenotes}
    \end{center}
\end{table}



% Chapter 3: 
\section{}

\subsection{}

\subsubsection{}

\subsubsection{}

\subsection{Differential Equation Model Based on System Dynamics}

The overall model can be expressed as:


\subsection{}

\subsubsection{Shannon Diversity Index}

\subsubsection{}


\subsection{}

\begin{table}[htbp]
    \begin{center}
        \caption{Specified or Initial Values}
        \begin{tabular}{c c c c}
            \toprule[2pt]
            \multicolumn{1}{m{2.5cm}}{\centering Parameters or Variables}
            &\multicolumn{1}{m{2.5cm}}{\centering Values}
            &\multicolumn{1}{m{2.5cm}}{\centering Parameters or Variables}
            &\multicolumn{1}{m{2.5cm}}{\centering Values}\\
            \midrule
            $r_P$ & 0.1 & $\beta_A$ & 0.1 \\
            $k_G$ & 1000 & $\mu_A$ & 0.05 \\
            $\alpha$ & 0.05 & $\beta_B$ & 0.1 \\
            $r_W$ & 0.5 & $\mu_B$ & 0.05 \\
            $\gamma$ & 0.05 & $P_0$ & 100 \\
            $\eta_I$ & 0.0001 & $W_0$ & 50 \\
            $\mu$ & 0.05 & $I_0$ & 30 \\
            $\mu'$ & 0.02 & $A_0$ & 1 \\
            $\eta$ & 0.05 & $B_0$ & 1 \\
            $\mu_r$ & 0.02 & $P_{r_0}$ & 20 \\
            \bottomrule[2pt]
        \end{tabular}
    \end{center}
\end{table}

Using Matlab, we calculated the temporal changes of these system dynamics variables
in Figure 6.

And we calculated the eigenvalue of Jacobian Matrix:
$$\lambda(J) = -0.0911, -0.1111, 0.2300, 0.0400$$





% Chapter 4:
\section{}
\subsection{}

\subsection{}

\begin{table}[htbp]
    \begin{center}
        \caption{Specified or Initial Values}
        \begin{tabular}{c c c c}
            \toprule[2pt]
            \multicolumn{1}{m{2.5cm}}{\centering Parameters or Variables}
            &\multicolumn{1}{m{2.5cm}}{\centering Values}
            &\multicolumn{1}{m{2.5cm}}{\centering Parameters or Variables}
            &\multicolumn{1}{m{2.5cm}}{\centering Values}\\
            \midrule
            $\eta_I$ & 0.005 & $\eta_{I_{l}}$ & 0.005 \\
            $\alpha_P$ & 1 & $\alpha_W$ & 15 \\
            $\gamma$ & 0.1 & $\mu_r$ & 1 \\
            $\mu_t$ & 1 & $\eta_P$ & 0.05 \\
            $\mu_{B_t}$ & 0.005  \\
            \bottomrule[2pt]
        \end{tabular}
        \begin{tablenotes}
            \footnotesize
            \item[*] *The values of the variables not listed 
            are the same as those in Table 2.
        \end{tablenotes}
    \end{center}
\end{table}



\begin{enumerate}
    \item \textbf{Enhancing Comprehensive Pest Control and Reducing Pesticide Dependence}

    \item \textbf{Enhancing Ecosystem Resilience}

    \item \textbf{Preventing Biodiversity Imbalance}
\end{enumerate}



% Chapter 5: 
\section{}
\subsection{Model Extension}



\subsection{}

\subsubsection{}

\subsubsection{}

\subsubsection{}


\begin{itemize}
    \item \textbf{Benefits:}
    \item[1] \textbf{Ecological Benefits:} 
    \item[2] \textbf{Economic Benefits:} 
    \item[3] \textbf{Sustainability Benefits:} 
    \item \textbf{Drawbacks}
    \item[1] \textbf{Management Complexity:} 
    \item[2] \textbf{Ecological Risks:} 
\end{itemize}


\subsubsection{}

\subsection{}



% Chapter 6: Strengths and Weaknesses
\section{Strengths and Weaknesses}
\subsection{Strengths}
\begin{itemize}
    \item \textbf{Logic:} 
    \item \textbf{Creativity:} 
    \item \textbf{Credibility:} 
    \item \textbf{Stability}: 
\end{itemize}


\subsection{Weaknesses}
\begin{itemize}
    \item 
    \item 
\end{itemize}




% References
\begin{thebibliography}{99}
    \bibitem{1} information on https://en.tutiempo.net/climate/cameroon.html.
\end{thebibliography}

\clearpage
\specialsectioning
\section{Letter}

\vspace{2cm}



\clearpage

\pagenumbering{gobble}
\section{Report on Use of AI}
\noindent
1.OpenAI \textit{ChatGPT}(https://chat.openai.com/) 

\label{1}
Query1:
Output:

\label{2}
Query2: 
Output: 

\label{3}
Query3: 
Output:
\newpage
\pagenumbering{arabic}
\setcounter{page}{26}


\end{document}